交通资产清查是支撑道路基础设施数字化管理、风险评估与精细化养护的重要基础工作。面向车载全景影像的交通资产清查需要同时解决“资产发现”和“状态理解”两类问题，但现有方法仍面临两方面瓶颈：复杂道路场景中的小目标、边缘目标与密集目标易发生漏检，影响资产发现完整性；异常状态样本稀缺且呈明显长尾分布，导致模型对低频高风险属性的识别能力不足。针对上述问题，本文以 WHU-Infra3D 数据集为基础，围绕交通资产清查中的“目标发现—属性理解—结构化记录”链条，研究主动感知检测与知识引导属性识别增强方法，旨在提升交通资产清查结果的完整性、准确性与工程可用性。

在目标检测方面，本文提出基于主动感知的多智能体三阶段交通资产检测框架。首先，对武汉与上海子集的等距柱状投影全景图进行四分透视投影切割，以缓解全景图局部几何拉伸对检测的干扰；随后，构建 Detection Agent 与 Crop Agent 协同的“全局初检—候选区域推荐—局部精检—结果融合”流程，将传统单次全图检测重构为面向潜在漏检区域的主动复查机制。在武汉训练集 6712 张透视图像和上海测试集 4700 张透视图像构成的跨城市实验中，所提框架在 RexOmni、DINO、Faster R-CNN、DEIM、Grounding DINO 和 YOLO26 六类异构检测器上均取得稳定增益，Macro-F1 与 Micro-F1 平均分别提升 8.18\% 和 6.41\%，并在小目标召回、密集区域补检和复杂场景目标恢复方面表现出更明显优势。

在属性识别方面，本文提出领域知识引导的长尾数据构建与增强方法。该方法将交通设施规范、运维经验与视觉先验转化为受约束的 Prompt 规则，在保持对象主体结构与场景上下文基本稳定的前提下，对锈蚀、风化老旧、结构损坏、井盖移位、垃圾满溢等关键异常状态进行定向补样，并构建“真实样本 + 知识引导生成样本”的混合训练数据。基于 WHU-Infra3D 构建的 7311 张真实局部样本，本文进一步生成 4197 条异常属性样本，并采用统一的多模态监督微调范式训练属性识别模型。实验结果表明，长尾增强模型在统一测试集上的样本级严格准确率、属性级准确率和异常属性准确率分别达到 88.78\%、96.12\% 和 93.72\%，较仅使用真实样本训练的模型分别提升 30.82、11.60 和 92.71 个百分点，正常—异常性能差距由 95.63\% 显著缩小至 6.08\%。

综上，本文从主动感知检测与知识引导长尾增强两个层面构建了面向交通资产清查的复合感知方法体系。研究结果表明，所提方法能够在保持工程可实现性的前提下，有效提升交通资产发现完整性、异常属性识别能力与结构化记录质量，可为交通资产数字化清查、智能巡检与养护决策提供方法支撑与数据基础。
